% ============================================================
%  Capítulo 1 — Introducción a la Inferencia Estadística
% ============================================================
\section{Introducción a la Inferencia Estadística}

\subsection{¿Qué es la inferencia estadística?}

La \textbf{inferencia estadística} es la rama de la estadística que se ocupa de
extraer conclusiones sobre una \emph{población} a partir de la información
contenida en una \emph{muestra}.

\begin{defi}[Población]
  Conjunto completo de individuos u observaciones sobre los cuales se desea
  obtener información.
\end{defi}

\begin{defi}[Muestra]
  Subconjunto de la población que se selecciona para su estudio.
\end{defi}

\begin{defi}[Parámetro]
  Característica numérica de la \emph{población} (generalmente desconocida),
  denotada habitualmente por letras griegas: $\theta$, $\mu$, $\sigma^2$, etc.
\end{defi}

\begin{defi}[Estadístico]
  Función de los datos de la muestra, $T = T(X_1, \ldots, X_n)$, que no depende
  de ningún parámetro desconocido.
\end{defi}

\subsection{Modelo estadístico}

Sea $X_1, X_2, \ldots, X_n \iid f(\cdot\,;\theta)$ una muestra aleatoria
simple (m.a.s.) de tamaño $n$ proveniente de una distribución con f.d.p.\ o
f.m.p.\ $f$, donde $\theta \in \Theta$ es el parámetro de interés.

\begin{teo}[Ley de los Grandes Números fuerte]
  Sea $X_1, X_2, \ldots$ una sucesión de variables aleatorias i.i.d.\ con
  $\E[X_1] = \mu < \infty$. Entonces
  \[
    \bar{X}_n = \frac{1}{n}\sum_{i=1}^{n} X_i \xrightarrow{\text{c.s.}} \mu
    \quad \text{cuando } n \to \infty.
  \]
\end{teo}

\begin{obs}
  La LGN justifica el uso de la media muestral $\bar{X}_n$ como estimador
  de $\mu = \E[X]$.
\end{obs}

\subsection{Tipos de inferencia}

\begin{enumerate}
  \item \textbf{Estimación puntual}: obtener un único valor $\hat{\theta}$ como
        estimación de $\theta$.
  \item \textbf{Estimación por intervalos}: construir un intervalo
        $[L, U]$ que contenga a $\theta$ con una probabilidad prefijada
        $(1-\alpha)$.
  \item \textbf{Pruebas de hipótesis}: decidir entre dos hipótesis sobre
        $\theta$ a partir de la evidencia muestral.
\end{enumerate}
